\section{Background}
\label{sec:background}
This section provides background information on video diffusion models~(Sec.~\ref{sec:video_diffusion}) and distribution matching distillation~(Sec.~\ref{sec:dmd}), which our method is built upon.

\subsection{Video Diffusion Models}
\label{sec:video_diffusion}
Diffusion models~\cite{sohl2015deep,ho2020denoising} generate samples from a data distribution $p(x_{0})$ by progressively denoising samples that are initially drawn from a Gaussian distribution $p(x_{T})$. 
They are trained to denoise samples created by adding random noise $\epsilon$ to the samples $x_0$ from the data distribution
\begin{equation}
    x_t = \alpha_t x_0 + \sigma_t\epsilon,\ \epsilon \sim \mathcal{N}(0, I)\,,
    \label{eq:diffusion_noise}
\end{equation}
where $\alpha_t, \sigma_T>0$ are scalars that jointly define the signal-to-noise ratio according to a specific noise schedule~\cite{karras2022elucidating,kingma2021variational,song2020score} at step $t$. 
The denoiser with parameter $\theta$ is typically trained to predict the noise~\cite{ho2020denoising}
\begin{equation}
    \mathcal{L}(\theta) = \mathbb{E}_{t, x_0, \epsilon}\left\|\epsilon_{\theta}(x_t, t) - \epsilon\right\|_2^2\,.
    \label{eq:denoising_loss}
\end{equation}
Alternative prediction targets include the clean image $x_0$~\cite{salimans2022progressive,karras2022elucidating} or a weighted combination of $x_0$ and $\epsilon$ known as v-prediction~\cite{salimans2022progressive}.
All prediction schemes are fundamentally related to the score function, which represents the gradient of the log probability of the distribution~\cite{song2020score,kingma2021variational}: \begin{equation} \label{eq
} s_\theta(x_t, t) = \nabla_{x_t} \log p(x_t) = - \frac{\epsilon_\theta(x_t, t)}{\sigma_t}. \end{equation} In the following sections, we simplify our notation by using the score function $s_\theta$ as a general representation of the diffusion model, while noting that it can be derived through reparameterization from a pretrained model of any prediction scheme.
At inference time, we start from full Gaussian noise $x_T$ and progressively apply the diffusion model to generate a sequence of increasingly cleaner samples. There are many possible sampling methods~\cite{song2020denoising,lu2022dpm,zhangfast,karras2022elucidating} to compute the sample $x_{t-1}$ in the next time step from the current one $x_t$ based on the predicted noise $\epsilon_{\theta}(x_t, t)$.

Diffusion models can be trained on either raw data~\cite{ho2020denoising,jabri2022scalable,karras2022elucidating} or on a lower-dimensional latent space obtained by a variational autoencoder~(VAE)~\cite{rombach2022high,peebles2023scalable,karras2024analyzing,yang2024cogvideox,opensora}.
The latter is often referred to as latent diffusion models~(LDMs) and has become the standard approach for modeling high-dimensional data such as videos~\cite{zhou2022magicvideo,opensora,hong2022cogvideo,blattmann2023align,yang2024cogvideox}. 
The autoencoder usually compresses both spatial and temporal dimensions of the video, making diffusion models easier to learn. The denoiser network %
in video diffusion models can be instantiated by different neural architectures, such as U-Net~\cite{ronneberger2015u,ho2022video,zhou2022magicvideo,chen2023videocrafter1} or Transformers~\cite{vaswani2017attention,hong2022cogvideo,yang2024cogvideox,videoworldsimulators2024}.

\subsection{Distribution Matching Distillation}
\label{sec:dmd}
Distribution matching distillation is a technique designed to distill a slow, multi-step teacher diffusion model into an efficient few-step student model~\cite{yin2024improved, yin2024onestep}. 
The core idea is to minimize the reverse KL divergence across randomly sampled timesteps $t$ between the smoothed data distribution $p_{\text{data}}(x_{t})$ and the student generator's output distribution $p_{\text{gen}}(x_{t})$. 
The gradient of the reverse KL can be approximated as the difference between two score functions:
\begin{align}
\label{eq:dmd}
\nabla_{\phi} \mathcal{L}_{\text{DMD}} & \triangleq \mathbb{E}_{t} \left( \nabla_{\phi} \text{KL} \left( p_{\text{gen}, t} \| p_{\text{data}, t} \right) \right) \notag \\
& \approx - \mathbb{E}_{t} \left( \int \left( s_{\text{data}} \left( \Psi \left( G_{\phi}(\epsilon), t \right), t \right) \right. \right. \\ 
& \qquad \left. \left. - s_{\text{gen},\xi} \left( \Psi \left( G_{\phi}(\epsilon), t \right), t \right) \right) \frac{d G_{\phi}(\epsilon)}{d \phi} \, d\epsilon \right), \notag 
\label{eq:dmd}
\end{align}
where $\Psi$ represents the forward diffusion process as defined in Eq.~\ref{eq:diffusion_noise}, $\epsilon$ is random Gaussian noise, $G_{\phi}$ is the generator parameterized by $\phi$, and $s_{\text{data}}$ and $s_{\text{gen},\xi}$ represent the score functions trained on the data and generator's output distribution, respectively, using a denoising loss~(Eq.~\ref{eq:denoising_loss}). 

During training, DMD~\cite{yin2024onestep} initializes both score functions from a pre-trained diffusion model. 
The score function of the data distribution is frozen, while the score function of the generator distribution is trained online using the generator's outputs.
Simultaneously, the generator receives gradients to align its output with the data distribution~(Eq.~\ref{eq:dmd}).  
DMD2~\cite{yin2024improved} extends this framework from single-step to multi-step generation by replacing the pure random noise input $\epsilon$ with a partially denoised intermediate image $x_t$.


